\documentclass[11pt]{article}
\usepackage{amsmath,amssymb}
\usepackage[a4paper,margin=1in]{geometry}
\usepackage{hyperref}

% Remove paragraph indentation
\setlength{\parindent}{0pt}
\setlength{\parskip}{6pt}

\begin{document}

\title{A Concise Derivation of UMAP\\
{\large Uniform Manifold Approximation and Projection}}
\author{Zhenwei Tang}
\date{}
\maketitle

% ============================================================
\section*{Goal}
% ============================================================

Given high-dimensional data points
\[
  x_i \in \mathbb{R}^p, \qquad i = 1, \dots, n,
\]
UMAP aims to find a low-dimensional embedding
\[
  y_i \in \mathbb{R}^d, \qquad d \ll p,
\]
such that the \emph{topological structure} of the data---captured by a weighted graph in the original space---is faithfully preserved in the embedding space.

The derivation proceeds in four stages:
\begin{enumerate}
  \item Construct a fuzzy topological representation of the high-dimensional data.
  \item Build an analogous fuzzy representation in the low-dimensional space.
  \item Define a cost function that measures the discrepancy between the two representations.
  \item Minimize the cost to obtain the embedding $\{y_i\}$.
\end{enumerate}

% ============================================================
\section*{Step 1: Local Metric Normalization}
% ============================================================

\textbf{Why this step?} Points in high-dimensional space are typically non-uniformly distributed. UMAP assumes the data lie on a Riemannian manifold that is \emph{locally} approximately uniform. To enforce this assumption, the distance metric is normalized locally for each point.

For each point $x_i$, let $\rho_i$ denote the distance to its nearest neighbor:
\[
  \rho_i = \min_{j \neq i} d(x_i, x_j),
\]
where $d(\cdot,\cdot)$ is the ambient Euclidean distance. This ensures the nearest neighbor of every point has effective distance zero, anchoring the local scale.

Let $k$ be the chosen number of neighbors. For each $x_i$, define the \emph{local kernel weight} from $x_i$ to $x_j$ as
\[
  v_{j \mid i} = \exp\!\left(-\frac{d(x_i, x_j) - \rho_i}{\sigma_i}\right),
\]
where $\sigma_i > 0$ is a bandwidth parameter chosen so that the \emph{effective number of neighbors} equals the perplexity target $\log_2 k$:
\[
  \sum_{j \neq i} v_{j \mid i} = \log_2 k.
\]

\textbf{Remark on $\sigma_i$.} This condition is solved numerically (e.g., by binary search) for each $i$ separately. The subtraction of $\rho_i$ in the exponent ensures $v_{j^* \mid i} = 1$ for the nearest neighbor $j^*$, preventing the kernel from collapsing in sparse regions.

% ============================================================
\section*{Step 2: Symmetrize into a Fuzzy Graph}
% ============================================================

\textbf{Why this step?} The weights $v_{j \mid i}$ are directed (asymmetric). To obtain an undirected graph suitable for a symmetric cost function, UMAP combines them using the \emph{probabilistic union}:
\[
  w_{ij} = v_{j \mid i} + v_{i \mid j} - v_{j \mid i} \cdot v_{i \mid j}.
\]

This formula is the fuzzy-set union: if $A$ and $B$ are two fuzzy membership values, then $A \cup B = A + B - A \cdot B$ (see Appendix~\ref{app:fuzzy}).

The result is a symmetric weighted graph $G = (V, E, w)$ with $w_{ij} \in (0, 1]$. High $w_{ij}$ means $x_i$ and $x_j$ are likely neighbors on the manifold.

% ============================================================
\section*{Step 3: Low-Dimensional Similarity}
% ============================================================

\textbf{Why this step?} We need an analogous similarity measure in the low-dimensional embedding space $\mathbb{R}^d$.

For the embedding coordinates $y_i, y_j \in \mathbb{R}^d$, UMAP defines the low-dimensional similarity as a \emph{Student-$t$} kernel (with one degree of freedom, matching t-SNE):
\[
  \hat{w}_{ij} = \left(1 + a \cdot \|y_i - y_j\|^{2b}\right)^{-1},
\]
where $a > 0$ and $b > 0$ are hyperparameters fit to approximate a smooth step function. Specifically, they are chosen so that
\[
  \hat{w}_{ij} \approx
  \begin{cases}
    1 & \text{if } \|y_i - y_j\| \leq \texttt{min\_dist}, \\
    \exp\!\left(-(\|y_i - y_j\| - \texttt{min\_dist})\right) & \text{otherwise},
  \end{cases}
\]
where \texttt{min\_dist} is a user-specified minimum distance in the embedding. This choice keeps nearby points tightly clustered while allowing global structure to spread out.

\textbf{Default values.} For the common setting \texttt{min\_dist} $= 0.1$, the fit yields $a \approx 1.929$ and $b \approx 0.791$.

% ============================================================
\section*{Step 4: The Cross-Entropy Cost Function}
% ============================================================

\textbf{Why this step?} UMAP treats $w_{ij}$ and $\hat{w}_{ij}$ as Bernoulli probabilities of an edge existing between $i$ and $j$ in the high- and low-dimensional graphs, respectively. The natural measure of discrepancy between two Bernoulli distributions is the \emph{binary cross-entropy} (see Appendix~\ref{app:crossentropy}).

The total cost is summed over all pairs:
\[
  \mathcal{L}(\{y_i\}) = \sum_{i} \sum_{j \neq i} \Big[
    w_{ij} \log \frac{w_{ij}}{\hat{w}_{ij}}
    + (1 - w_{ij}) \log \frac{1 - w_{ij}}{1 - \hat{w}_{ij}}
  \Big].
\]

Dropping terms that do not depend on $\{y_i\}$ (i.e., terms involving $w_{ij} \log w_{ij}$ and $(1-w_{ij})\log(1-w_{ij})$), the \emph{effective} minimization objective reduces to:
\[
  \boxed{
    \mathcal{L}(\{y_i\}) = -\sum_{(i,j) \in E} w_{ij} \log \hat{w}_{ij}
    - \sum_{(i,j) \notin E} (1 - w_{ij}) \log (1 - \hat{w}_{ij}).
  }
\]

\textbf{Interpretation.}
\begin{itemize}
  \item The \textbf{first term} (attractive): for pairs that are close in the original space ($w_{ij} \approx 1$), penalize large $\|y_i - y_j\|$.
  \item The \textbf{second term} (repulsive): for pairs that are distant in the original space ($w_{ij} \approx 0$), penalize small $\|y_i - y_j\|$.
\end{itemize}

% ============================================================
\section*{Step 5: Optimization via Stochastic Gradient Descent}
% ============================================================

\textbf{Why this step?} The full double sum over all $O(n^2)$ pairs is computationally prohibitive. UMAP uses \emph{Stochastic Gradient Descent (SGD)} with \emph{negative sampling} to approximate the gradients efficiently.

\textbf{Gradient with respect to $y_i$.} We compute $\nabla_{y_i} \mathcal{L}$.

Denote $r_{ij} = \|y_i - y_j\|^2$ and $\hat{w}_{ij} = (1 + a \cdot r_{ij}^b)^{-1}$.

\textbf{Attractive gradient} (for edge $(i, j) \in E$, contribution from the first term):

Applying the chain rule:
\[
  \frac{\partial}{\partial y_i}\bigl[-w_{ij} \log \hat{w}_{ij}\bigr]
  = -w_{ij} \cdot \frac{1}{\hat{w}_{ij}} \cdot \frac{\partial \hat{w}_{ij}}{\partial r_{ij}} \cdot \frac{\partial r_{ij}}{\partial y_i}.
\]

Computing each factor:
\[
  \frac{\partial \hat{w}_{ij}}{\partial r_{ij}} = -a b \, r_{ij}^{b-1} \hat{w}_{ij}^2,
  \qquad
  \frac{\partial r_{ij}}{\partial y_i} = 2(y_i - y_j).
\]

Substituting:
\[
  \nabla_{y_i}^{\mathrm{attr}} = \frac{2 a b \, w_{ij} \, r_{ij}^{b-1}}{1 + a \, r_{ij}^b} \cdot (y_i - y_j).
\]

\textbf{Repulsive gradient} (for non-edge $(i, j) \notin E$, contribution from the second term):

\[
  \frac{\partial}{\partial y_i}\bigl[-(1-w_{ij})\log(1-\hat{w}_{ij})\bigr]
  = (1-w_{ij}) \cdot \frac{1}{1-\hat{w}_{ij}} \cdot \frac{\partial \hat{w}_{ij}}{\partial r_{ij}} \cdot \frac{\partial r_{ij}}{\partial y_i}.
\]

Since for non-edges $w_{ij} \approx 0$:
\[
  \nabla_{y_i}^{\mathrm{rep}} \approx \frac{-2 b \, \hat{w}_{ij}}{(1-\hat{w}_{ij})(1 + a\,r_{ij}^b)} \cdot r_{ij}^{b-1} \cdot (y_i - y_j).
\]

\textbf{Update rule.} At each SGD step, for a sampled positive edge $(i,j)$ and $m$ sampled negative pairs $(i, k_1), \ldots, (i, k_m)$:
\[
  y_i \leftarrow y_i - \eta \left(\nabla_{y_i}^{\mathrm{attr}} + \sum_{\ell=1}^{m} \nabla_{y_i}^{\mathrm{rep},\ell}\right),
\]
where $\eta$ is a decaying learning rate. This reduces per-step cost from $O(n^2)$ to $O(m)$.

% ============================================================
\section*{Conclusion}
% ============================================================

UMAP starts by modeling the high-dimensional data as a fuzzy topological graph with edge weights
\[
  w_{ij} = v_{j \mid i} + v_{i \mid j} - v_{j \mid i} \cdot v_{i \mid j},
\]
derived from locally normalized distances. It then seeks low-dimensional coordinates $\{y_i\}$ such that the embedding's similarity structure
\[
  \hat{w}_{ij} = \left(1 + a\|y_i - y_j\|^{2b}\right)^{-1}
\]
matches the original graph as closely as possible. The match is quantified by the cross-entropy cost
\[
  \mathcal{L} = -\sum_{(i,j)\in E} w_{ij} \log \hat{w}_{ij}
  - \sum_{(i,j)\notin E} (1-w_{ij})\log(1-\hat{w}_{ij}),
\]
which is minimized via SGD with negative sampling.

\begin{center}
\fbox{\parbox{0.88\textwidth}{\centering
UMAP embeds data by minimizing the cross-entropy between a fuzzy high-dimensional graph\\
and an analogous low-dimensional similarity structure,\\
preserving local topology via locally normalized kernels.
}}
\end{center}

% ============================================================
\appendix
\section*{Appendices}
% ============================================================

\subsection*{Appendix A: Fuzzy Set Union}\label{app:fuzzy}

In fuzzy set theory, membership values are real numbers in $[0,1]$. For two fuzzy memberships $A, B \in [0,1]$, the \emph{fuzzy union} (probabilistic OR) is defined as:
\[
  A \cup B = A + B - A \cdot B.
\]
This is the inclusion-exclusion formula for two independent Bernoulli events with probabilities $A$ and $B$:
\[
  P(A \text{ or } B) = P(A) + P(B) - P(A)P(B).
\]
Setting $A = v_{j\mid i}$ and $B = v_{i\mid j}$ gives the symmetrized edge weight $w_{ij}$.

\subsection*{Appendix B: Binary Cross-Entropy}\label{app:crossentropy}

For two Bernoulli distributions with parameters $p$ and $q$, the \emph{cross-entropy} of $q$ relative to $p$ is:
\[
  H(p, q) = -p \log q - (1-p)\log(1-q).
\]
It measures how well the distribution $q$ approximates $p$: $H(p,q) \geq H(p,p)$, with equality iff $p = q$ (Gibbs' inequality). In UMAP, $p = w_{ij}$ (high-dimensional) and $q = \hat{w}_{ij}$ (low-dimensional). Summing over all pairs gives the total cost $\mathcal{L}$.

\subsection*{Appendix C: Riemannian Manifold Assumption}\label{app:manifold}

UMAP assumes the data $\{x_i\}$ lie on (or near) a smooth Riemannian manifold $\mathcal{M} \subset \mathbb{R}^p$. A Riemannian manifold is a smooth $d$-dimensional manifold equipped with a smoothly varying inner product (metric tensor) $g_x$ on each tangent space $T_x\mathcal{M}$.

Crucially, UMAP assumes the manifold is \emph{locally uniform}: with respect to the Riemannian metric, data points are uniformly distributed. Since the true metric is unknown, UMAP approximates it by normalizing distances locally using $\rho_i$ and $\sigma_i$, effectively learning a point-adaptive metric.

\subsection*{Appendix D: Stochastic Gradient Descent (SGD)}\label{app:sgd}

Standard gradient descent minimizes a loss $\mathcal{L}(\theta) = \sum_i \ell_i(\theta)$ by iterating:
\[
  \theta \leftarrow \theta - \eta \nabla_\theta \mathcal{L}(\theta).
\]
When the sum over all pairs is too expensive, \emph{Stochastic Gradient Descent} approximates the full gradient by sampling a mini-batch $\mathcal{B}$ at each step:
\[
  \theta \leftarrow \theta - \eta \, \nabla_\theta \sum_{(i,j) \in \mathcal{B}} \ell_{ij}(\theta).
\]
In UMAP, each positive edge is sampled once per epoch, and $m$ negative pairs are sampled uniformly at random, giving an unbiased stochastic estimate of $\nabla \mathcal{L}$ per step. The learning rate $\eta$ is typically annealed as $\eta(t) = \eta_0 (1 - t/T)$ over $T$ total epochs.

\subsection*{Appendix E: Negative Sampling}\label{app:negsampling}

Computing the repulsive term exactly requires summing over all $O(n^2)$ non-edges. Negative sampling approximates this by, for each positive edge $(i,j)$, drawing $m$ random non-neighbors $k_1, \ldots, k_m$ uniformly and computing their repulsive gradients only. With $m$ typically set to 5, this reduces cost from $O(n^2)$ to $O(n m)$ per epoch. The approach originates from Word2Vec and is valid because non-edges are the vast majority of pairs, so uniform sampling provides a low-variance gradient estimate.

\end{document}
